\subsection{Importância das Variáveis e dos Sensores}

% Nomenclatura provisória solicitada: o modelo arbóreo executado é ExtraTreesClassifier.
Para comparar a contribuição das entradas nos dois classificadores, foi
realizada uma análise adicional com os mesmos conjuntos de treinamento e teste,
separados por coleta. Em cada classificador foram ajustadas duas versões: uma
com os seis sinais MQ brutos (MQ2, MQ3, MQ7, MQ8, MQ135 e MQ138) e outra com
esses sinais acrescidos do índice de umidade do solo, da temperatura e da
pressão interna. Os sinais MQ brutos foram mantidos nas duas versões porque os
MQ corrigidos utilizados na análise anterior já incorporam uma compensação
estatística baseada nas variáveis ambientais. Assim, os resultados desta
ablação não substituem os percentuais apresentados nas tabelas anteriores.
O modelo arbóreo utilizou 900 árvores, seleção de atributos pela raiz
quadrada do número de entradas e mínimo de dez amostras por folha. A X-TCN
foi treinada por 60 épocas com janelas de 128 leituras e passo de 64; a
semente aleatória foi 42 nos dois casos.

O teste incluiu 22.115 leituras de 12 coletas, das quais 16.692 leituras de
oito coletas foram classificadas operacionalmente como solo seco e 5.423
leituras de quatro coletas como solo úmido. Cada modelo foi treinado uma única
vez com todas as coletas de treinamento; a separação por condição do solo foi
aplicada apenas à avaliação. Na X-TCN, a probabilidade atribuída a cada
leitura corresponde à média das probabilidades das janelas temporais que a
contêm. A Tabela~\ref{tab:ablacao_ambiente_vacuo} apresenta os resultados
obtidos por leitura.

\begin{table}[htbp]
    \centering
    \caption{Desempenho dos classificadores com e sem variáveis ambientais nas coletas a vácuo.}
    \label{tab:ablacao_ambiente_vacuo}
    \scriptsize
    \begin{tabular}{llrrrr}
        \hline
        \textbf{Modelo} & \textbf{Entradas} &
        \textbf{Acurácia total (\%)} &
        \textbf{Bal. total (\%)} &
        \textbf{Bal. seco (\%)} &
        \textbf{Bal. úmido (\%)} \\
        \hline
        Random Forest & MQ              & 89,86 & 89,90 & 99,93 & 38,56 \\
        Random Forest & MQ + ambiente   & 87,93 & 87,89 & 99,92 & 33,30 \\
        X-TCN         & MQ              & 80,06 & 79,77 & 91,50 & 26,07 \\
        X-TCN         & MQ + ambiente   & 89,41 & 89,54 & 95,12 & 46,15 \\
        \hline
    \end{tabular}

    \vspace{0.1cm}
    {\footnotesize Bal. = acurácia balanceada. Fonte: Elaborado pelo autor (2026).}
\end{table}

Para medir a importância das entradas nos dois modelos, foi adotada a
\emph{importância por permutação} (\emph{permutation feature importance}),
conforme o princípio de perturbar uma variável por vez e medir a alteração do
desempenho do modelo \cite{fisher2019modelreliance}. Em cada repetição, os
valores de uma coluna foram embaralhados entre as leituras do teste completo,
mantendo-se o modelo treinado e as demais colunas. A importância corresponde à
diferença entre a acurácia balanceada original e a obtida após a permutação,
expressa em pontos percentuais. Foram realizadas três permutações por variável
em cada cenário. O mesmo procedimento e a mesma métrica foram usados para o
Random Forest e para a X-TCN; nesta última, as janelas foram reconstruídas a
partir das leituras perturbadas antes da predição. A
Tabela~\ref{tab:importancia_variaveis_vacuo} apresenta a média das três quedas.

\begin{table}[htbp]
    \centering
    \caption{Importância por permutação no teste completo: queda da acurácia balanceada, em pontos percentuais.}
    \label{tab:importancia_variaveis_vacuo}
    \scriptsize
    \begin{tabular}{lrrrr}
        \hline
        \textbf{Variável} &
        \textbf{RF: MQ} &
        \textbf{X-TCN: MQ} &
        \textbf{RF: MQ + amb.} &
        \textbf{X-TCN: MQ + amb.} \\
        \hline
        Pressão interna & --    & --    & 7,13  & 23,49 \\
        Umidade do solo & --    & --    & -0,37 & 8,12  \\
        Temperatura     & --    & --    & -2,60 & -1,01 \\
        MQ2             & 2,40  & -3,68 & -0,24 & 17,64 \\
        MQ3             & 3,84  & -8,92 & 1,13  & -3,77 \\
        MQ7             & 3,88  & -1,14 & 1,63  & 3,07  \\
        MQ8             & 9,28  & -5,11 & 1,96  & 5,25  \\
        MQ135           & 9,40  & 8,20  & 0,57  & 3,77  \\
        MQ138           & 0,17  & -5,13 & 1,43  & 10,59 \\
        \hline
    \end{tabular}

    \vspace{0.1cm}
    {\footnotesize -- = variável ausente do modelo. Fonte: Elaborado pelo autor (2026).}
\end{table}

Com as variáveis ambientais, a pressão interna produziu a maior queda de
acurácia balanceada nos dois modelos: 7,13 pontos percentuais no Random Forest
e 23,49 na X-TCN. Na X-TCN, também se destacaram MQ2 (17,64), MQ138 (10,59)
e o índice de umidade do solo (8,12 pontos percentuais). Sem as entradas
ambientais, MQ135 apresentou a maior queda nos dois classificadores. Essas
diferenças mostram que a posição de um sensor no ranking depende das demais
entradas oferecidas ao modelo. Valores negativos indicam que, nessa avaliação,
o embaralhamento da variável aumentou a acurácia balanceada; não representam
uma influência física negativa do sensor sobre a presença de nematoides.

A inclusão das variáveis ambientais elevou a acurácia balanceada total da
X-TCN de 79,77\% para 89,54\%, mas reduziu a do Random Forest de 89,90\% para
87,89\%. No recorte úmido, a X-TCN passou de 26,07\% para 46,15\% e o Random
Forest passou de 38,56\% para 33,30\%. Portanto, o ganho da X-TCN com as
informações ambientais não eliminou a dificuldade de classificação nessa
condição. A distinção entre solo seco e úmido foi feita por um índice
operacional do sensor de solo, sem medição calibrada da água em diferentes
profundidades. Como o teste úmido contém apenas quatro coletas e a condição de
umidade também varia com o dia de coleta, esses resultados não isolam um efeito
causal da água sobre os gases ou sobre os nematoides. A importância preditiva
da pressão também pode refletir características da câmara e da operação do
sistema de vácuo.
